\section{Organizzazione e Metodologia del lavoro}

Nel corso del progetto, l'approccio organizzativo del gruppo si è evoluto progressivamente. Nelle fasi iniziali le attività sono state ripartite in modo orizzontale (``per compiti''), data l'assenza di una netta differenziazione dei ruoli o di un'assegnazione basata su competenze pregresse. Successivamente, la suddivisione è avvenuta sulla base delle attitudini individuali, definendo ruoli operativi ben precisi: \textit{Marco Piovesan} si è concentrato sul perfezionamento di alcune sezioni del \texttt{D1} e sull'integrazione tra backend e frontend; \textit{Michele Porcellato} ha curato lo sviluppo del codice del backend, la redazione del \texttt{D2} e assistito \textit{Marco} durante l'integrazione tra frontend e backend; \textit{Alessio Pesaresi} si è occupato della modellazione dei diagrammi del \texttt{D3} e della stesura del business plan.

\paragraph{Incontri}
    Per monitorare lo stato di avanzamento delle attività e affrontare tempestivamente eventuali criticità, il gruppo ha istituito momenti di confronto a cadenza almeno settimanale. Tali riunioni sono state determinanti per concordare variazioni architetturali o funzionali, assegnare i task individuali e stabilire gli obiettivi da raggiungere per la settimana successiva. Inoltre, la convivenza di tutti e tre i componenti nella medesima struttura ha agevolato costanti allineamenti informali e una rapida definizione delle scadenze interne.

\paragraph{Strumenti utilizzati}
    La redazione della documentazione è stata realizzata in \LaTeX{} tramite un progetto condiviso su Overleaf, mentre per il controllo versione del codice sorgente è stato utilizzato un repository GitHub. Il monitoraggio del tempo impiegato (dettagliato nel capitolo~\ref{caricoLavoro}) è stato costantemente aggiornato mediante un documento condiviso, successivamente evolutosi in una bacheca per l'annotazione di idee e la gestione delle attività da completare.